% ECONOMICS_PROBLEM: {"id": "C73-30", "title": "Weak master-equation selection beyond finite-state potentials", "jel_code": "C73", "source_id": "OP-0216"}
% FIRST_POSTED: 2026-10-10
% ATTEMPT_STATUS: solved
\documentclass[11pt]{article}
\usepackage[a4paper,margin=25mm]{geometry}
\usepackage{amsmath,amssymb,amsthm,mathtools,microtype,titling}
\usepackage[hidelinks]{hyperref}
\setlength{\droptitle}{-8mm}
\newtheorem{theorem}{Theorem}
\newtheorem{lemma}{Lemma}
\newtheorem{corollary}{Corollary}
\newcommand{\E}{\mathbb E}
\newcommand{\R}{\mathbb R}
\newcommand{\eps}{\varepsilon}
\newcommand{\cA}{\mathcal A}
\newcommand{\cD}{\mathcal D}
\DeclareMathOperator{\Var}{Var}
\DeclareMathOperator{\Cov}{Cov}
\title{Shock selection in a nonpotential\protect\\ finite-state master equation}
\author{Ji Ho Bae}
\date{10 October 2026}
\hypersetup{pdftitle={Shock selection in a nonpotential finite-state master equation},pdfauthor={Ji Ho Bae}}
\begin{document}
\maketitle
\begin{abstract}
We give an explicit family of smooth nonpotential costs for which the uncorrected finite-state master equation has a unique selected weak feedback field. Every admissible vanishing boundary repulsion gives the same local $L^1$ limit, without a rate condition between the repulsion and the common noise. The limit has an intrinsic scalar entropy characterization. The family includes interior shocks: its permitted time horizon is independent of front steepness, and a scaling extends the result to every finite horizon. The proof uses localization of the passive values and a Cole--Hopf comparison construction that remains valid through shock formation.
\end{abstract}

\section{The problem and a class of data}
Fix $d\ge3$ and the simplex $\Delta_d=\{m_i\ge0:\sum_i m_i=1\}$. Write
\[
 H_i(U)=\frac12\sum_{j\ne i}(U_i-U_j)_+^2,\qquad
 b_i(m,U)=\sum_{j\ne i}\{m_j(U_j-U_i)_+-m_i(U_i-U_j)_+\}.
\]
The uncorrected regularization considered here is
\begin{align}\label{master}
0={}&\partial_tU_i+\varepsilon\sum_{j,l}(m_j\delta_{jl}-m_jm_l)\partial^2_{m_jm_l}U_i
 +(b+R^\phi)\cdot DU_i+2\varepsilon(e_i-m)\cdot DU_i-H_i(U)\notag\\
&+f_i(m)+\sum_{j\ne i}\phi_\varepsilon(m_j)(U_j-U_i),\qquad U_i(T,m)=g_i(m),
\end{align}
where $R_i^\phi=\phi_\varepsilon(m_i)-m_i\sum_j\phi_\varepsilon(m_j)$ and derivatives are tangent to the simplex. Following the selection formulation in catalogue C73-30, let $\mathcal A(f,g,T)$ consist of smooth nonnegative families $\phi_\varepsilon$, converging to zero uniformly on every compact subset of $(0,1]$, for which \eqref{master} has a unique classical solution continuous on the closed cylinder, with bounded value and first derivatives and locally bounded second derivatives in the interior. No bound on these derivatives uniform in $\varepsilon$ is imposed. The feedback field is $(U_i-U_d)_{i<d}$.

The problem asks for nonpotential assumptions and an intrinsic admissibility condition giving existence of selected limits, uniqueness independent of the admissible regularization, and agreement with classical solutions. The potential theory and this extension are discussed in \cite[\S4.3--4.4]{CD}. We address the sufficient-condition question with a family that undergoes loss of classical regularity.

Put $q=(d-1)/2$, $p=m_1$, and $r=(r_2,\ldots,r_{d-1})$, where $r_j=m_j-m_d$. Thus
\[
 m_d=\frac{1-p-\sum_{j=2}^{d-1}r_j}{d-1},\qquad m_j=m_d+r_j.
\]
Choose fixed smooth cutoffs $\chi:\mathbb R\to[0,1]$ and $\beta:\mathbb R^{d-2}\to[0,1]$ such that
\[
 \chi=0\ \text{on }(-\infty,1/128],\quad \chi=1\ \text{on }[1/64,\infty),
\]
and, with $R_d=1/(16d)$, $\beta=1$ on $\{|r|_\infty\le R_d/4\}$ and $\beta=0$ on $\{|r|_\infty\ge R_d/2\}$. Set $\alpha=1/16$, $y=\sqrt p$, $\theta=y-1/2-\alpha r_2$, and, for $k\ge1$, define
\begin{equation}\label{data}
 f=0,\qquad g=(G_k,0,\ldots,0),\qquad
 G_k(p,r)=\frac1q\left[24(1-\chi(p))+\chi(p)\frac{2-\beta(r)\tanh(k\theta)}{\sqrt p}\right].
\end{equation}
The second term is defined as zero where $\chi=0$. These data are smooth on a neighborhood of $\Delta_d$, and $1/q\le G_k\le34/q$ there on the simplex.

For each admissible $r$, let $0<y<L(r)$ be the interval corresponding to the simplex interior, where
\[
 L(r)^2=1-\sum_{j=2}^{d-1}r_j+(d-1)\min(0,r_2,\ldots,r_{d-1})>0.
\]
Define $v_0(y,r)=qyG_k(y^2,r)$, extend it by zero to $y<0$, and extend it constantly to the right of $L(r)$. The construction gives $0\le v_0\le3$. Let $v(s,y,r)$ be the bounded nonnegative entropy solution of
\begin{equation}\label{burgers}
 v_s+\partial_y(v^2/2)=0,\qquad v(0,y,r)=v_0(y,r).
\end{equation}
The value on $0<y<L(r)$ is independent of the chosen nonnegative right extension. Indeed, its Hopf--Lax minimizers lie in $[0,y]$: the initial primitive is nondecreasing, and the quadratic penalty increases for a footpoint to the right of $y$. The zero extension also gives the left trace zero. We use the usual bounded entropy uniqueness and vanishing-viscosity theorem for this scalar Cauchy problem \cite{Kru}.

\begin{theorem}\label{main}
For each $d\ge3$ and the fixed cutoffs above, there is a $T_0>0$, independent of $k\ge1$, with the following properties. For $0<T\le T_0$ and the data \eqref{data}, the class $\mathcal A(f,g,T)$ is nonempty. For every $\phi\in\mathcal A(f,g,T)$,
\begin{equation}\label{convergence}
 U^{\varepsilon,\phi}\longrightarrow
 U^*(t,m)=\left(\frac{v(T-t,\sqrt p,r)}{q\sqrt p},0,\ldots,0\right)
 \quad\text{in }L^1_{\rm loc}((0,T)\times\operatorname{int}\Delta_d).
\end{equation}
Consequently all selected feedback fields coincide, and any two admissible families have feedback distance tending to zero on every interior compact set.

Among bounded feedback fields, an intrinsic admissibility condition is that the passive components vanish and that their remaining nonnegative component, multiplied by $q\sqrt p$ and expressed in time $s=T-t$, satisfies \eqref{burgers} in the entropy sense with its initial trace. This condition determines the selected field uniquely and agrees with the classical master field whenever the latter exists.

The data are nonpotential even modulo a common component. If $kT>1$, the selected field has an interior shock for $1/k<s\le T$. Thus the assertion includes a nonclassical regime, while retaining every family in $\mathcal A$.
\end{theorem}

\section{Bounds that do not use derivatives uniform in the noise}
The following estimates also justify the reduction to one active value. All constants may depend on the fixed data and on the compact set, but not on $\varepsilon$ or on derivatives of $\phi_\varepsilon$.

\begin{lemma}\label{localization}
Suppose $f=0$, $g=(G,0,\ldots,0)$ and $0<\gamma\le G\le M$. Every admissible solution satisfies $0\le U_i\le M$. For $D_\eta=\{m_i\ge\eta\ \forall i\}$, put
\[
 C=(d-1)M+1,\quad \delta=\tfrac12\eta e^{-CT},\quad
 a_\varepsilon=\sup_{[\delta,1]}\phi_\varepsilon,\quad
 e_\varepsilon=d\exp\left(-\frac{\eta^2}{4\varepsilon T e^{2CT}}\right).
\]
For all sufficiently small $\varepsilon$, on $[0,T]\times D_\eta$,
\begin{equation}\label{passive}
 0\le U_i\le M(Ta_\varepsilon+e_\varepsilon)\quad(i>1),\qquad
 U_1\ge\gamma e^{-[qM+(d-1)a_\varepsilon]T}(1-e_\varepsilon).
\end{equation}
\end{lemma}
\begin{proof}
First fix $\varepsilon$. The ordinary maximum principle extends to the degenerate boundary as follows. Add $hB_a(m)=h\sum_jm_j^a$, $0<a<1$, to a component being maximized. Bounded first derivatives force the maximum into the interior: adding a small empty coordinate gains order $h\rho^a$, whereas the other changes are $O(\rho)$. At this maximum the Hamiltonian and value-switching terms have the appropriate nonpositive signs. The drift without $R^\phi$ has coordinate lower bounds $-C_\varepsilon m_j$, since the solution is bounded for this fixed $\varepsilon$. The repulsion contributes
\[
 -hR^\phi\cdot DB_a=-ha\sum_j\phi_\varepsilon(m_j)\{m_j^{a-1}-B_a(m)\}.
\]
Because $B_a\le d^{1-a}$, each summand with $m_j\le1/d$ is nonpositive. The remaining coordinates lie in $[1/d,1]$, where $\phi_\varepsilon$ is bounded. Thus this contribution is $O_\varepsilon(h)$ as well, regardless of the size of the repulsion close to zero.

If $m_l\ge1/d$, the tangent first-order condition bounds
$ha m_j^{a-1}$ by $L_\varepsilon+ha d^{1-a}$, where $L_\varepsilon$ bounds the first derivatives. The diffusion error is at most
\[
 \varepsilon d(1-a)(L_\varepsilon+ha d^{1-a}),
\]
and the drift error is $O_\varepsilon(h)$. A strict time barrier, followed by $h\downarrow0$ and then $a\uparrow1$, gives $U_i\le M$. Applied to $-U$ it gives $U_i\ge0$; at a minimum component $H_i=0$. No uniform bound on $L_\varepsilon$ was used.

Stop an interior diffusion, with the Kimura covariance and drift
$b(m,U)+R^\phi+2\varepsilon(e_I-m)$, on leaving $D_\delta$ or reaching $T$. The label $I$ may switch in any adapted way. Once $(d-1)a_\varepsilon+2\varepsilon\le1$, each drift coordinate is at least $-Cm_j$. After multiplication by $e^{C(s-t)}$, its finite-variation part is nondecreasing. Its martingale part has quadratic variation at most $\varepsilon T e^{2CT}/2$, because the coordinate covariance rate is $2\varepsilon m_j(1-m_j)\le\varepsilon/2$. An exit requires one such martingale to fall by at least $\eta/2$. The exponential martingale inequality and a union bound give exit probability at most $e_\varepsilon$.

For a passive component, write
\[
 -H_i(U)=\sum_{j\ne i}c_{ij}(U_j-U_i),\qquad c_{ij}=\tfrac12(U_i-U_j)_+.
\]
Its $j=1$ term is nonpositive and can be dropped for an upper comparison. The repulsion term involving $U_1$ is at most $M\phi_\varepsilon(m_1)$. Keep the nonnegative switching rates between passive labels. The stopped cooperative-system representation has terminal value zero, exit value at most $M$, and source at most $Ma_\varepsilon$. This proves the first estimate in \eqref{passive}. All coefficients are bounded on the stopped interior domain for fixed $\varepsilon$.

Finally $H_1\le qU_1^2\le qMU_1$, and its repulsion value term is at least $-(d-1)a_\varepsilon U_1$. The stopped scalar comparison with killing rate $qM+(d-1)a_\varepsilon$, terminal value $\gamma$ and exit value zero gives the second estimate.
\end{proof}

We also record a local consistency consequence used at the boundary of the shock region. If a bounded classical deterministic master field $V$ is defined on a protected spatial neighborhood for the whole time interval and has bounded first and second derivatives there, then $U^{\varepsilon,\phi}\to V$ uniformly on any smaller neighborhood from which a path with bounded drift cannot leave the protected region by time $T$. One also stops inside $D_\delta$ as above. Here is the full error estimate behind this statement. Linearizing the positive parts gives a cooperative switching matrix, while $(b(U)-b(V))DV$ is a bounded, possibly signed, zeroth-order matrix. Replacing its off-diagonal entries by their absolute values bounds both signs of the error by a cooperative system with a bounded potential. The residual obtained by inserting $V$ is at most $C(\varepsilon+a_\varepsilon)$. Therefore
\begin{equation}\label{localclassical}
 |U^{\varepsilon,\phi}-V|\le C\{\varepsilon+a_\varepsilon+\mathbb P(\text{exit before }T)\}.
\end{equation}
The first exit through $D_\delta$ is controlled by the preceding proof. For the other boundary, a fixed margin larger than the maximal drift displacement leaves a fixed martingale displacement necessary for exit; its probability tends to zero by the $L^2$ maximal inequality. This proves \eqref{localclassical} without assuming bounds on derivatives of $U^{\varepsilon,\phi}$ uniform in $\varepsilon$.

\section{Reduction and the regular collars}
Write $u=U_1$. On an interior compact cylinder, \eqref{passive} makes all passive components tend uniformly to zero and leaves $u$ uniformly positive. Elementary algebra in the rates and Hamiltonians gives
\begin{equation}\label{scalarperturb}
 u_s+q\partial_p(pu^2)
 =\varepsilon\mathcal L u+E_\varepsilon\cdot D_{p,r}u+F_\varepsilon,
 \qquad \|E_\varepsilon\|_\infty+\|F_\varepsilon\|_\infty\longrightarrow0.
\end{equation}
Here $\mathcal L$ is the scalar Kimura operator including the extra drift $2(e_1-m)\cdot D$, and the errors contain no derivatives of passive values. Indeed,
$b_1=-2qpu+O(\max_{i>1}U_i)$, $b_j-b_d=O(\max_{i>1}U_i)$, $H_1=qu^2+O(\max_{i>1}U_i)$, and the repulsion terms have uniformly small coefficients on the compact set. Treat these error functions, obtained from the actual solution, as fixed coefficients in \eqref{scalarperturb}.

Under $V=qyu$, the unperturbed part of \eqref{scalarperturb} becomes $V_s+VV_y$. In $(y,r)$ the second-order coefficients are
\begin{equation}\label{coefficients}
 A_{yy}=\frac{1-y^2}{4},\qquad A_{yr_j}=-\frac{yr_j}{2},\qquad
 A_{r_jr_l}=m_j\delta_{jl}+m_d-r_jr_l.
\end{equation}
The first-order coefficients are $(1-y^2)/(4y)$ and $-r_j$, and the zeroth-order coefficient is $-(1-y^2)/(4y^2)$. After the fixed change $\theta=y-1/2-\alpha r_2$, \eqref{scalarperturb} is exactly
\begin{equation}\label{Veq}
 V_s+VV_\theta=\varepsilon\mathcal D V+B_\varepsilon\cdot D_{\theta,r}V+C_\varepsilon,
 \qquad \|B_\varepsilon\|_\infty+\|C_\varepsilon\|_\infty\longrightarrow0.
\end{equation}
On interior compact sets, $\mathcal D$ has fixed smooth coefficients and a positive definite second-order matrix. Denote its $\partial_{\theta\theta}$ coefficient by $a(\theta,r)>0$. The conversion of the errors is harmless because $y$ is bounded away from zero.

Consider the rectangle
\[
 \Omega=\{(\theta,r):|\theta|<1/8,\ |r|_\infty<R_d\}.
\]
Its closure and a slightly larger rectangle lie in the simplex interior, and $\chi=1$ there. Thus the initial value for \eqref{Veq} is exactly $2-\beta(r)\tanh(k\theta)$ on this neighborhood.

We now choose $T_0$, before choosing $k$. All derivatives through any fixed order of the initial data are bounded uniformly in $k$ outside the smaller strip
\[
 \mathcal S=\{|\theta|\le1/16,\ |r|_\infty\le R_d/2\}.
\]
This follows from $\sup_{k\ge1}k^j\operatorname{sech}^2(k/16)<\infty$; in the remaining region $\beta$ vanishes. The cutoff near $p=0$ is fixed and stays away from this strip. Insert a fixed closed rectangle $\mathcal S_1$ strictly between $\mathcal S$ and $\Omega$, for instance with widths $3/32$ and $3R_d/4$ in the displayed coordinates. Their Euclidean separations in the simplex are positive.

Let $C_1$ bound the negative part of $\partial_yv_0$ outside $\mathcal S$, uniformly in $k$. Choose $T_0$ so that $T_0C_1<1/2$, and so that the following displacements are less than half the respective fixed separations: the deterministic characteristic displacement between $\mathcal S$ and $\mathcal S_1$, and the maximal stopped population drift displacement between $\mathcal S_1$ and the collars of $\partial\Omega$. Such bounds are uniform in $k$: $v_0\le3$, $U_i\le M=34/q$, and on a stopped interior domain $\phi_\varepsilon\le1$, $2\varepsilon\le1$ eventually, giving a drift bound depending only on $d,M$. Finally take $3T_0$ smaller than the width of the larger front collar, and $2T_0<1/32$. These are finitely many positive requirements independent of $k$.

For the deterministic ansatz $(u,0,\ldots,0)$ the equations reduce exactly to \eqref{burgers}, with $r$ constant along characteristics. Its characteristic map is $y=z+sv_0(z,r)$. A minimizing footpoint differs from $y$ by at most $3s$. Hence, outside $\mathcal S_1$, all possible footpoints avoid $\mathcal S$; on their interval the action has second derivative at least $1/s-C_1>0$. There is one smooth branch, with characteristic Jacobian at least $1/2$. This constructs a classical field with bounded derivatives on the boundary collars and locally outside $\mathcal S_1$.

For \eqref{localclassical}, stop on approaching $\mathcal S_1$ as well as on leaving $D_\delta$. Starting on a boundary collar or on a compact set outside $\Omega$, the choice of $T_0$ leaves a fixed martingale displacement necessary to approach $\mathcal S_1$. Its probability tends to zero. Thus \eqref{localclassical} applies and the actual boundary values converge uniformly to the classical field. No classical field is assumed across the shock.

\section{A uniform Cole--Hopf estimate}
Let $w(s,x;\nu,\lambda)$ solve viscous Burgers on $\mathbb R$, with viscosity $\nu>0$ and initial value $2-\lambda\tanh(kx)$, where $0\le\lambda\le1$. Write
\begin{equation}\label{action}
 S(z)=2z-\frac\lambda k\log\cosh(kz)+\frac{(x-z)^2}{2s},\qquad
 w=\frac{x-\mathbb E_\nu z}{s},\qquad
 d\mu_\nu(z)=\frac{e^{-S(z)/(2\nu)}dz}{\int e^{-S/(2\nu)}}.
\end{equation}
This is the Cole--Hopf formula; differentiating the heat kernel verifies the equation and initial value. In this section expectations and covariances use $\mu_\nu$. All estimates are uniform for $x$ in a fixed compact interval, $s$ in a fixed bounded interval, and $\lambda\in[0,1]$. The parameter $k$ is fixed.

\begin{lemma}\label{CH}
Put $Q=-w_x\ge0$. There is a modulus $\omega(\nu)\to0$ such that
\begin{align}
 |w|&\le3,\qquad |\nu w_{xx}|+|w_\lambda|\le C(1+Q),\label{CH1}\\
 |\nu w_\nu|+|\nu^2w_{x\nu}|+|\nu^3w_{\nu\nu}|
 +|\nu^2w_{\nu\lambda}|+|\nu w_{x\lambda}|+|\nu w_{\lambda\lambda}|
 &\le\omega(\nu)(1+Q).\label{CH2}
\end{align}
The estimates hold from $s=0$ by continuous extension.
\end{lemma}
\begin{proof}
The maximum principle gives $1\le w\le3$ and $w_x\le0$. For $s\ge s_0>0$, put $Z=z-\mathbb E z$, $A=S-\mathbb E S$, and $R=\partial_\lambda S=-k^{-1}\log\cosh(kz)$. Differentiation under the integral gives
\begin{align*}
 \operatorname{Var}(z)&=2\nu s^2(Q+1/s),&
 w_{xx}&=-\frac{\mathbb E Z^3}{4\nu^2s^3},\\
 w_\nu&=-\frac{\mathbb E ZA}{2\nu^2s},&
 w_{x\nu}&=\frac{\operatorname{Var}(z)}{2\nu^2s^2}-\frac{\mathbb E Z^2A}{4\nu^3s^2},\\
 w_{\nu\nu}&=\frac{\mathbb E ZA}{\nu^3s}-\frac{\mathbb E ZA^2}{4\nu^4s},&
 w_\lambda&=\frac{\operatorname{Cov}(z,R)}{2\nu s},\\
 w_{x\lambda}&=\frac{\mathbb E Z^2(R-\mathbb E R)}{4\nu^2s^2},&
 w_{\lambda\lambda}&=-\frac{\mathbb E Z(R-\mathbb E R)^2}{4\nu^2s},\\
 w_{\nu\lambda}&=-\frac{\operatorname{Cov}(z,R)}{2\nu^2s}
 +\frac{\mathbb E Z(R-\mathbb E R)A}{4\nu^3s}.&&
\end{align*}

Here are uniform moment estimates, including the possibly degenerate transition at shock formation. On a bounded $z$ interval the derivatives $S'',S''',S''''$ do not vanish simultaneously: if $\lambda=0$ then $S''=1/s$; otherwise $S'''$ vanishes only at zero, where $S''''=2\lambda k^3>0$. Outside a sufficiently large interval $S''$ is uniformly positive. A finite cover consequently gives
\[
 |\{z:S(z)-\min S\le h\}|\le C h^{1/4}\quad(0<h\le1).
\]
For completeness, the elementary sublevel estimate used here follows by choosing $j+1$ separated points of a set on which $|F|\le h$. Divided differences bound $\inf|F^{(j)}|$ by $C_jh/|\{|F|\le h\}|^j$. Apply this on each covering interval where one derivative of order $j\le4$ has a fixed nonzero sign. The quadratic tails of $S$ handle larger sublevels.

At a global minimizer, the uniform upper bound on $S''$ gives a Gaussian lower bound $c\sqrt\nu$ for the partition integral after subtraction of $\min S$. Integrating the sublevel estimate, or summing energy strips of width $\nu$, gives
\begin{equation}\label{actionmom}
\mathbb E(S-\min S)^j+\mathbb E|A|^j\le C\nu^{j-1/4}\qquad(1\le j\le4).
\end{equation}
Moreover $\mathbb E z\in[x-3s,x-s]$, by integration by parts in \eqref{action}. Gaussian far tails, together with $\operatorname{Var}(z)\ge2\nu s$, imply
\begin{equation}\label{zmom}
 \mathbb E|Z|^3+\mathbb E Z^4\le C\operatorname{Var}(z).
\end{equation}
Indeed, on a fixed bounded interval the higher powers are bounded by a constant times $Z^2$; the complementary moments are exponentially small compared with $\nu$. Equations \eqref{actionmom}--\eqref{zmom} and Cauchy--Schwarz prove the viscosity-derivative bounds in \eqref{CH2}, for example with factor $C\nu^{3/8}$, and the bound on $\nu w_{xx}$. Since $R$ is $1$-Lipschitz, $|\operatorname{Cov}(z,R)|\le\operatorname{Var}(z)$ and $\mathbb E[Z^2(R-\mathbb E R)^2]\le C\operatorname{Var}(z)$. This also proves the bounds on $w_\lambda$ and $\nu^2w_{\nu\lambda}$.

The remaining two derivatives need a sharper observation. Set $h=x-2s$. Apart from an additive constant,
\[
 S(z)=z^2/(2s)-hz/s-(\lambda/k)\log\cosh(kz).
\]
If $h\ne0$, there is a unique global minimizer, of the sign of $h$. To see this for $h>0$, compare $z$ with $-z$ and then use that $S''$ is increasing on the positive half-line and $S'(0)<0$. If $h=0$, the minimizer is zero or a symmetric pair. Hence the even function $R$ has the same value at every global minimizer, for every parameter.

Compactness now implies that the oscillation of $R$ on the sublevel $S-\min S\le a$ tends uniformly to zero as $a\downarrow0$. Outside a fixed such sublevel the Gibbs moments are exponentially small. On the sublevel use this small oscillation for one factor and the Lipschitz covariance bound for the other factors. Since $\operatorname{Var}(z)\ge c\nu$, it follows that
\[
 \mathbb E[Z^2(R-\mathbb E R)]=o(\operatorname{Var}(z)),\qquad
 \mathbb E[Z(R-\mathbb E R)^2]=o(\operatorname{Var}(z)),
\]
uniformly in the parameters. For the second estimate one uses
$\mathbb E|Z||R-\mathbb E R|\le\operatorname{Var}(z)$; centering by $\mathbb E R$ changes the sublevel oscillation only by the same vanishing modulus and an exponentially small tail. The displayed derivative identities prove the desired bounds on $\nu w_{x\lambda}$ and $\nu w_{\lambda\lambda}$.

Finally choose $s_0\le1/(2k)$. Before this time the differentiated Burgers equation gives $-2k\le w_x\le0$. Successive spatial derivatives satisfy linear parabolic equations with bounded coefficients and already bounded lower derivatives. Differentiating in $\nu$ or $\lambda$ gives the same type of equations, with initial data zero or derivatives of $2-\lambda\tanh(kx)$; the source for $w_\nu$, for example, is $w_{xx}$. Induction through the spatial order four and parameter order two gives bounds independent of $\nu$ on this initial interval. Each scaled derivative in \eqref{CH2} then tends to zero there as well. This completes the estimates through $s=0$.
\end{proof}

\section{Comparison through the shock}
On a slightly larger rectangle than $\Omega$, set
\[
 W_\varepsilon(s,\theta,r)=w(s,\theta;\varepsilon a(\theta,r),\beta(r)).
\]
The chain rule and Lemma~\ref{CH} give, with a modulus $\eta_\varepsilon\to0$,
\begin{align}
 1-W_\theta&\asymp1+Q,\qquad W_\theta\le\eta_\varepsilon,
 \qquad |DW|+\varepsilon|D^2W|\le C(1-W_\theta),\label{weighted}\\
 |W_s+WW_\theta-\varepsilon\mathcal D W|&\le\eta_\varepsilon(1-W_\theta).\label{residual}
\end{align}
We spell out why the potentially large derivatives cause no omitted term. Since $\beta$ is independent of $\theta$, $W_\theta=w_x+O(\nu w_\nu)$ with $\nu=\varepsilon a$. The main normal diffusion cancels $w_s+ww_x=\nu w_{xx}$. The remaining normal terms involve $\nu w_\nu$, $\nu^2w_{x\nu}$ and $\nu^3w_{\nu\nu}$. Mixed and tangential diffusion adds $\nu w_{x\lambda}$, $\nu w_{\lambda\lambda}$ and $\nu^2w_{\nu\lambda}$, all small in the weighted sense. Terms with a single $w_\lambda$ carry a factor $\varepsilon$. The fixed lower-order coefficients are bounded. This proves \eqref{weighted}--\eqref{residual}.

The boundary collars constructed earlier have only regular characteristics. The Cole--Hopf measures there concentrate uniformly on their unique minimizers, so $W_\varepsilon$ converges uniformly to the same classical field as the actual $V_\varepsilon$. Let $b_\varepsilon\to0$ be the sum of these two uniform errors on a collar of $\partial\Omega$, and let $L$ bound the classical field's $\theta$ derivative there. Let $\Delta_\varepsilon\to0$ bound the two error coefficients in \eqref{Veq} on the rectangle.

Choose nonnegative functions $a_0(s),h(s)$ by
\begin{equation}\label{odebarrier}
 K=C(h+a_0+\eta_\varepsilon+\Delta_\varepsilon),\quad
 h'=a_0+K,\quad a_0'=Lh'+K,\quad
 h(0)=0,\quad a_0(0)=b_\varepsilon,
\end{equation}
where $C$ is a sufficiently large fixed constant. They tend uniformly to zero. For small $\varepsilon$ the shifts stay inside the collar. The functions
\[
 W^+(s,\theta,r)=W_\varepsilon(s,\theta-h(s),r)+a_0(s),\qquad
 W^-(s,\theta,r)=W_\varepsilon(s,\theta+h(s),r)-a_0(s)
\]
are respectively a supersolution and a subsolution of \eqref{Veq}. Indeed, translation of the fixed smooth coefficients changes the residual by at most $Ch(1-W_\theta)$, using the Hessian bound in \eqref{weighted}. The other errors, including the zeroth-order term on $a_0$, are at most $K(1-W_\theta)$. The upper residual is therefore at least
\[
 (-h'+a_0)W_\theta+a_0'-K(1-W_\theta)=a_0'-K\ge0;
\]
the lower residual is at most $K-a_0'\le0$.

The initial inequalities hold because $W_\varepsilon(0)=V_\varepsilon(0)$ on the rectangle. On its boundary, \eqref{odebarrier} gives $a_0-Lh\ge b_\varepsilon$, which proves the boundary inequalities. At each fixed $\varepsilon$ the scalar equation is uniformly parabolic on this interior rectangle, so the comparison principle yields $W^-\le V_\varepsilon\le W^+$. Only boundedness, not derivatives, of the error coefficients was used.

For each $r$, the reference tends to the Burgers entropy solution, except possibly on the shock plane $\theta=2s$ when $\beta(r)ks>1$. This is also the solution defined from the full data \eqref{data}: both sets of minimizing footpoints lie within distance $3T_0$ of the observation point, where the initial functions coincide; their primitives differ only by an additive constant. Boundedness gives convergence in local $L^1$. Also $W_\theta\le\eta_\varepsilon$ and $1\le W\le3$ give a uniform bound on its total variation on each fixed $\theta$ interval. Its small normal shifts therefore tend to zero in $L^1$, uniformly in the other variables. The comparison just proved gives $V_\varepsilon\to v$ in $L^1$ on the rectangle. Outside it, \eqref{localclassical} gives local uniform convergence. Lemma~\ref{localization} removes the passive components. This proves \eqref{convergence} for every admissible family.

\section{Admissibility, nonemptiness and scope}
The change $(p,r)\mapsto(y,r)$ has a smooth nonzero Jacobian on interior compact sets, so the convergence above is precisely in the simplex measure required in the question. It implies compactness and pairwise independence of the regularizations. The entropy condition \eqref{burgers}, imposed for almost every parameter $r$, is intrinsic to the limiting deterministic equation and has no reference to $\varepsilon$ or $\phi$. All selected feedback fields are bounded by the value estimate.

For clarity, entropy uniqueness also applies on the parameter-dependent interval used here. The local Kato inequality for two nonnegative entropy solutions has flux
$\tfrac12(v+\widetilde v)|v-\widetilde v|\ge0$. Integrate with a cutoff that increases near zero and decreases near $L(r)$. Bounded feedback gives $v,\widetilde v=O(y)$ near zero, so the left cutoff error vanishes. The right cutoff contribution is nonpositive and can be discarded. The initial $L^1$ distance is zero; hence the solutions agree. This is the usual entropy contraction argument \cite[Theorem~1]{Kru}, with the inflow and outflow terms made explicit. It proves uniqueness under the stated intrinsic condition.

If a classical full master solution exists, follow one of its characteristics. The value vector solves $\dot u_i=H_i(u)$ backwards from $(G_k,0,\ldots,0)$. Uniqueness of this locally Lipschitz ODE keeps the passive values zero and the active value positive. The full classical solution therefore has the scalar reduction \eqref{burgers}, and agrees with its entropy solution. This verifies classical consistency without assuming the ansatz for an arbitrary candidate classical field.

The admissible class is nonempty. Bayraktar, Cecchin, Cohen and Delarue \cite[Theorem~3.4, equation~(2.19)]{BCCD} provide a classical Wright--Fisher solution for smooth costs when the repulsion is sufficiently large near zero and vanishes beyond an arbitrarily small boundary layer. Their noise parameter is $\sqrt{2\varepsilon}$ in our normalization. Their first derivatives extend continuously to the boundary. Uniqueness extends to the class used here: subtract two fixed-$\varepsilon$ solutions, linearize, and apply the same global barrier maximum principle used in Lemma~\ref{localization}, now with a bounded signed zeroth-order matrix and zero terminal data. Taking the layer width to zero gives an admissible family. Only sufficiently small $\varepsilon$ matters for the limit; if a family is required at larger $\varepsilon$ too, the substitution $U=cV$, $\tau=ct$, $\phi=c\widetilde\phi$, with $c>2\varepsilon$, reduces it to the same theorem.

At $p=1/4,r=0$, the intrinsic terminal covector is $G_k\,dm_1$. Varying $m_2$ with $m_d$ compensating gives
\[
 \partial_{m_2}G_k=\frac{4\alpha k}{q}\ne0,
 \qquad \partial_{m_1}(g_2-g_d)=0.
\]
It is not a gradient, even after a common component is added to the costs. On the core where $\beta=1$, the initial Burgers derivative is $-k$ at $\theta=0$. Characteristics first cross at $s=1/k$; for later times the symmetric decreasing front has its entropy shock at $\theta=2s$. The construction of $T_0$ did not depend on $k$, so this shock occurs strictly inside the permitted cylinder whenever $kT>1$.

This supplies the requested nonpotential assumptions, an intrinsic weak admissibility condition, and all three selection conclusions, including a genuine post-classical regime. It is a sufficient-condition result for the stated data family, not a claim that arbitrary nonpotential costs always have a unique vanishing-noise limit. No cost correction, prescribed rate for $\phi_\varepsilon$, or uniform bound on noisy derivatives is part of the theorem.

\begin{corollary}[Every finite horizon]\label{allhorizons}
For any $d\ge3$ and $T>0$, choose $\sigma>0$ with $\sigma T\le T_0$ and use $f=0$, $g=(\sigma G_k,0,\ldots,0)$. All the existence, independence, intrinsic admissibility and classical-consistency conclusions of Theorem~\ref{main} hold. The entropy initial value is $\sigma v_0$, and an interior shock occurs when $\sigma kT>1$.
\end{corollary}
\begin{proof}
Set $\widehat U(\tau,m)=\sigma^{-1}U(\tau/\sigma,m)$, $\widehat\varepsilon=\varepsilon/\sigma$, and
$\widehat\phi_{\widehat\varepsilon}=\sigma^{-1}\phi_{\sigma\widehat\varepsilon}$. The homogeneities of $H$ and $b$ show directly that \eqref{master} becomes the same equation on the horizon $\sigma T$, with terminal value $(G_k,0,\ldots,0)$. This transformation is a bijection of the respective admissible regularization classes; their regularity, uniqueness and local vanishing properties are preserved. Theorem~\ref{main} applies, and scalar entropy solutions scale as $v^\sigma(s,y,r)=\sigma v(\sigma s,y,r)$. One may take $\sigma=\min(1,T_0/(2T))$ and then choose $k>1/(\sigma T)$ to obtain a shock on any prescribed horizon.
\end{proof}

\begin{thebibliography}{9}
\small
\bibitem{CD} P. Cardaliaguet and F. Delarue.
\emph{Selected topics in mean field games}.
Proceedings of the International Congress of Mathematicians 2022, vol.~5, 3660--3703.
DOI: \href{https://doi.org/10.4171/ICM2022/17}{10.4171/ICM2022/17}.
\bibitem{BCCD} E. Bayraktar, A. Cecchin, A. Cohen and F. Delarue.
\emph{Finite state mean field games with Wright--Fisher common noise}.
\href{https://arxiv.org/pdf/1912.06701v2}{arXiv:1912.06701v2}, 10 January 2021, Theorem~3.4 and \S3.2.2.
\bibitem{Kru} S. N. Kruzhkov.
\emph{First order quasilinear equations in several independent variables}.
Mathematics of the USSR-Sbornik 10(2) (1970), 217--243.
DOI: \href{https://doi.org/10.1070/SM1970v010n02ABEH002156}{10.1070/SM1970v010n02ABEH002156}.
\end{thebibliography}
\smallskip\noindent\footnotesize
Prepared with Codex assistance; the verification record describes the mathematical and source checks.
\normalsize
\section*{Submission scope and verification}
This answers the sufficient-condition formulation of C73-30: the explicit smooth nonpotential family in every $d\geq3$ and every finite horizon $T>0$ satisfies all three requested conclusions. The arbitrary-horizon scaling in Corollary~\ref{allhorizons} preserves the entire admissible class. Every family in the original admissible boundary-repulsion class converges to the same weak feedback field, with intrinsic entropy admissibility and classical consistency. The permitted horizon is independent of the front steepness, so the family includes interior shocks. No cost correction, rate comparison between repulsion and noise, or uniform bound on noisy derivatives is assumed.

This does not claim convergence for arbitrary nonpotential costs; the original question asks to find assumptions. The earlier smooth-regime conditional results are not reused as a whole solution.

Author: Ji Ho Bae. Prepared and checked through GPT Codex; the exact serving revision was not exposed. The authoring assistant performed the analytic checks. No new independent-agent audit, external peer review, or Lean verification is claimed.
\begin{itemize}
\item \href{https://github.com/jbaelaw/nonpotential-master-shock-selection/blob/4c9b0ae26e615c098b9044dd5ee0eca480222e1e/paper/Proof.pdf}{Complete eight-page proof PDF} and \href{https://github.com/jbaelaw/nonpotential-master-shock-selection/blob/4c9b0ae26e615c098b9044dd5ee0eca480222e1e/paper/Proof.tex}{TeX source}.
\item \href{https://github.com/jbaelaw/nonpotential-master-shock-selection/blob/4c9b0ae26e615c098b9044dd5ee0eca480222e1e/verification/Review.md}{Mathematical review} and \href{https://github.com/jbaelaw/nonpotential-master-shock-selection/blob/4c9b0ae26e615c098b9044dd5ee0eca480222e1e/verification/SourceAudit.md}{primary-source audit}.
\item \href{https://github.com/jbaelaw/nonpotential-master-shock-selection/blob/4c9b0ae26e615c098b9044dd5ee0eca480222e1e/verification/Verification.json}{Verification record} and \href{https://github.com/jbaelaw/nonpotential-master-shock-selection/blob/4c9b0ae26e615c098b9044dd5ee0eca480222e1e/provenance/Process.md}{process and attribution}.
\end{itemize}
The original statement, rank, and earlier attempts are preserved. Submitted for expert review; catalogue-maintainer acceptance is not asserted.


\end{document}
